\chapter{Исследование существующих подходов}

\startrelatedwork

Настоящая глава содержит обзор существующих способов хранения
полуструктурированных данных в СУБД. Рассматриваются четыре
принципиально различных подхода: хранение JSON как неделимого
значения, модель Entity-Attribute-Value, колоночное развёртывание
путей, и гибридный тип \texttt{JSON} в ClickHouse.

Каждый из подходов анализируется по трём характеристикам:
эффективность ввода-вывода при запросе к отдельному полю,
локальность кэша при обработке большого числа строк, и стоимость
эволюции схемы при появлении нового поля.

\section{Хранение JSON как неделимого значения}

Исторически первый и до сих пор наиболее распространённый способ
хранения JSON в СУБД~--- рассматривать документ целиком как
значение одной ячейки особого типа. Документ сериализуется в
текстовом виде (JSON) или в бинарном (BSON в MongoDB, JSONB в
PostgreSQL).

Текстовый JSON (MySQL до 5.7, ранние версии Oracle) хранит JSON как
тип CLOB без внутреннего индексирования: при каждом обращении к полю
происходит полный разбор строки парсером.

Каноническим представителем документной модели на сегодняшний день
является MongoDB~\cite{mongodb-bson}, разработавшая бинарный формат
BSON и фактически закрепившая в индустрии саму концепцию
<<документоориентированной СУБД>>. В BSON документ предварён
таблицей смещений внутренних элементов, каждое скалярное значение
хранится в нативном бинарном формате, типы значений
самоописательны. Близкое по идее представление JSONB в PostgreSQL
дополнительно упорядочивает ключи объектов для ускорения бинарного
поиска. Оба подхода позволяют обращаться к отдельному полю за время
$O(\log n)$ без полного разбора, однако два фундаментальных
недостатка сохраняются:

\begin{enumerate}
    \item \textbf{Полное чтение документа с диска.} Страница диска в
    СУБД типично составляет 4--16~КБ, а один JSON-документ нередко
    имеет размер от единиц до десятков килобайт. Для запроса
    \texttt{SELECT doc.user.id FROM t WHERE doc.kind = 'commit'}
    система должна считать с диска каждый документ целиком.
    \item \textbf{Отсутствие колоночного доступа.} Несколько
    последовательных строк таблицы содержат свои документы,
    разнесённые по страницам; чтение одного и того же поля для
    миллиона строк приводит к миллиону разрозненных обращений к
    памяти.
\end{enumerate}

Эмпирически на аналитических нагрузках над JSON-данными
документная модель и JSONB демонстрируют производительность на
один-два порядка хуже колоночных СУБД на том же объёме данных.
Индексы (составные, функциональные, в случае MongoDB~--- в том
числе wildcard) частично ускоряют выборку по конкретному пути, но
не решают фундаментальной проблемы ввода-вывода при последующем
чтении других полей и не ускоряют агрегатные запросы,
затрагивающие всю таблицу. Именно эти ограничения сделали
документную модель плохо приспособленной к OLAP-нагрузкам и
послужили мотивацией для появления специализированных колоночных
форматов хранения JSON.

\section{Модель Entity-Attribute-Value}

EAV~--- классическая реляционная техника моделирования данных с
переменной структурой, описанная в литературе ещё в 1970--1980-е
годы в контексте медицинских и биологических баз
данных~\cite{stead1983,nadkarni2007}. Идея состоит в декомпозиции
гетерогенных записей в универсальную тройку:
\begin{verbatim}
CREATE TABLE attr (
  entity_id BIGINT,
  attr_name TEXT,
  attr_value TEXT
);
\end{verbatim}

Документ \texttt{\{"a":~1, "b":~"x"\}} превращается в две строки:
\texttt{(1, 'a', '1')} и \texttt{(1, 'b', 'x')}. Запрос
\texttt{SELECT a FROM t} принимает вид
\texttt{SELECT attr\_value FROM attr WHERE attr\_name~=~'a'}.

Сильные стороны: схема полностью динамическая, отсутствие
NULL-значений, простая модель мутаций. Слабые стороны, многократно
описанные в литературе~\cite{dinu2007}:

\begin{itemize}
    \item \textbf{Потеря типизации.} Все значения приводятся к
    строке. Агрегация числовых полей требует преобразования типов
    на каждой строке.
    \item \textbf{Деградация запросов с несколькими полями.} Запрос
    \texttt{SELECT a, b, c FROM t} становится трёхкратным
    self-join либо группировкой с условными агрегатами. Сложность
    превышает колоночный вариант на порядок.
    \item \textbf{Иерархия.} Для вложенных JSON-документов
    возникает выбор между двумя не-эквивалентными решениями:
    через \texttt{parent\_id} (рекурсивный обход) и через
    разворот путей (flattened~paths, \cite{beyer2005,florescu1999}).
\end{itemize}

Обе вариации EAV страдают от общего недостатка: колоночный доступ к
одному полю невозможен без предикатного сканирования по
\texttt{attr\_name}. Как следствие, аналитические нагрузки
деградируют многократно. Это явление формулируется в работах по
декомпозиции хранилищ как <<vertical~fragmentation>> проблема
групповой обработки атрибутов~\cite{copeland1985}.

\section{Колонка на каждый путь JSON}

Альтернативный подход, к которому пришли современные аналитические
движки (Snowflake VARIANT~\cite{snowflake-variant}, ClickHouse JSON
в ранних версиях, ряд коммерческих систем)~--- разворачивать
JSON-документ в колоночное представление, где каждый уникальный
путь становится отдельной физической колонкой.

Документ \texttt{\{"a":~1, "b":~"x", "c":\{"d":~2\}\}} при вставке
порождает три колонки: \texttt{a}~(INT), \texttt{b}~(STRING),
\texttt{c.d}~(INT). Отсутствующие в конкретном документе поля
записываются как NULL.

Сильные стороны: колоночный доступ в полной мере, полная типизация,
применимость всех стандартных оптимизаций (сжатие, zone~maps,
векторизация).

Слабые стороны:

\begin{itemize}
    \item \textbf{Рост числа физических колонок.} В наборе JSONBench
    из публичного потока сети Bluesky (около 40~млн событий) при
    полной развёртке обнаруживается более 10~тысяч уникальных
    JSON-путей, при этом 99\% событий используют лишь около 50
    путей. Классическая СУБД с дисковыми страницами фиксированного
    размера не справляется с таблицей в 10~000 колонок.
    \item \textbf{Разрастание NULL-значений.} При $N~=~10^7$ строк и
    $5\cdot 10^3$ редких колонок только на битмаски валидности
    требуется порядка 6~ГБ.
    \item \textbf{Стоимость ALTER TABLE ADD COLUMN.} Каждый вызов
    добавления колонки требует $O(N)$ работы по инициализации
    значений по умолчанию для всей истории таблицы.
\end{itemize}

Теоретические основы колоночного хранения описаны в
C-Store~\cite{cstore}, MonetDB/X100~\cite{monetdb-x100},
Vectorwise~\cite{vectorwise}. Во всех этих работах рассматривается
случай фиксированной известной заранее схемы; случай динамической
схемы с произвольным ростом числа колонок не был предметом
систематического анализа до недавнего времени.

\section{Тип данных JSON в ClickHouse}

Наиболее продвинутым промышленным решением в области колоночного
хранения полуструктурированных данных на сегодня является новый
тип данных \texttt{JSON} в ClickHouse, введённый в версии~24.8 и
стабилизированный в~24.10~\cite{clickhouse-json}. Схема хранения
устроена следующим образом.

\textbf{Частично развёрнутая схема.} Тип декларируется как
\begin{verbatim}
JSON(
  max_dynamic_paths = N,
  field_1 Type_1,
  field_2 Type_2,
  ...
)
\end{verbatim}

Часть путей объявляется явно (с зафиксированным типом); они
становятся полноценными колонками основной таблицы со всеми
свойствами колоночного хранения (сжатие, zone~maps,
специализированные кодеки). Остальные пути помещаются в общее
хранилище \texttt{shared~data}.

\textbf{Shared~data.} Редкие пути хранятся внутри строки таблицы
как пара вложенных массивов: \texttt{\{"paths":[path\_1, \ldots],
"values":[value\_1, \ldots]\}}. По сути, это EAV-представление,
упакованное внутрь одной ячейки.

\textbf{Динамическое продвижение.} При вставке новых данных
система отслеживает частоту появления путей; пути, частота которых
превышает заданный порог \texttt{max\_dynamic\_paths} (по умолчанию
1024), автоматически <<продвигаются>> в полноценные колонки на
следующем слиянии.

Сильные стороны: автоматическое распознавание часто встречающихся
путей, эффективный колоночный доступ к популярным полям, компактное
хранение длинного хвоста.

Ограничения:
\begin{itemize}
    \item Ограниченный набор операций над \texttt{shared~data}:
    агрегации по путям оттуда существенно медленнее, чем по
    полноценным колонкам; индексы к ним не применяются.
    \item Отсутствие кросс-документной колоночной локальности для
    \texttt{shared~data}: чтобы узнать значение редкого поля в
    100~строках, ClickHouse должен разобрать 100~массивов.
    \item Стоимость слияний: перенос пути из \texttt{shared~data} в
    отдельную колонку требует переписывания всех кусков хранения
    (parts), затронутых полем.
\end{itemize}

Важно, что в текущей реализации ClickHouse редкие пути хранятся
внутри каждой строки, что означает дублирование ключей пути в
каждой строке (с применением словарной компрессии, но всё же в
логически миллионной кратности для часто встречающегося, но не
продвинутого пути).

\section{Выводы главы}

Проведённый анализ позволяет сформулировать требования к
разрабатываемому решению:
\begin{enumerate}
    \item Часто встречающиеся поля должны храниться как полноценные
    колонки~--- необходимое условие для достижения
    производительности, сопоставимой с ClickHouse, на типичных
    аналитических запросах.
    \item Редко встречающиеся поля должны храниться отдельно от
    основной таблицы: прямолинейное раздувание схемы неприемлемо;
    размещение редких значений внутри каждой строки сохраняет
    проблему чтения лишних данных при колоночных запросах.
    \item Должна существовать процедура автоматической миграции
    полей между режимами по мере роста их частоты, без
    необходимости ручного вмешательства в схему.
    \item Для разреженных полей нужно сохранить возможность
    применения стандартных операций (предикаты, индексы,
    агрегации)~--- желательно без специализированного кода, за
    счёт переиспользования существующих механизмов реляционного
    движка.
\end{enumerate}

Совокупно эти требования приводят к решению, описанному в
следующей главе: хранение каждого редкого пути в отдельной
side-таблице формата \texttt{(\_id, value)}, гибридное сканирование
одним специализированным оператором с подстановкой значений по
\texttt{\_id}, и явное разделение полей на закреплённые (pinned)
и разреженные (sparse) на этапе \texttt{CREATE TABLE}.

\finishrelatedwork

\chapterconclusion

Изложенные в настоящей главе подходы демонстрируют, что
существующие решения распределены между двумя крайностями:
документные системы теряют преимущества колоночного доступа,
классическое колоночное хранилище не справляется с произвольно
расширяющейся схемой. Гибридные подходы ClickHouse и Snowflake
частично закрывают этот разрыв, однако вносят собственные
ограничения. Предлагаемое в настоящей работе решение основано на
переиспользовании полноценных реляционных таблиц в качестве
хранилища для разреженных полей, что позволяет избежать
специализированного формата с ограниченной функциональностью.
